%% Supplementary material for the revised manuscript. All tables and numbers are generated.
\documentclass[10pt]{article}
\usepackage[T1]{fontenc}
\usepackage{times}
\usepackage[margin=1.8cm]{geometry}
\usepackage{booktabs,graphicx,xcolor,amsmath,amssymb,longtable,url,algorithm,algpseudocode}
\makeatletter
\input{generated/numbers.tex}
\input{generated/numbers_extra.tex}
\input{generated/numbers_expansion.tex}
\newcommand{\R}[1]{\ifcsname R@#1\endcsname\csname R@#1\endcsname\else\textbf{\textcolor{red}{[#1]}}\fi}
\makeatother
\newcommand{\sg}{\operatorname{sg}}
\renewcommand{\thetable}{S\arabic{table}}
\renewcommand{\thefigure}{S\arabic{figure}}

\title{Supplementary material: When does boundary distillation help lightweight medical image segmentation?}
\date{}
\begin{document}
\maketitle
\noindent Citations refer to the reference list of the main article. ``Original'' and ``retrained'' are the model sets defined in the main text. Arm codes used in some tables: M0 U-Net, M6 BDS-Lite, M1 U-Net with surface loss only, M2 U-Net with boundary decoder only, M7 U-Net with boundary decoder and surface loss, M3--M5 BDS-Lite without surface loss, distillation, or gate; GSL U-Net+GSL, C1 U-Net+Boundary Loss, C3 Attention U-Net, C4 LB-UNet. SESOI labels: ``within SESOI'' means the 95\% interval lies inside $\pm$SESOI, ``beyond'' means it lies wholly outside, and ``inconclusive'' means neither.

\section*{S1. Boundary-aware methods: contributions and limitations}
\begin{table}[ht]\centering\small
\resizebox{\linewidth}{!}{%
\begin{tabular}{p{2.6cm}p{1.8cm}p{3.8cm}p{1.6cm}p{2.6cm}p{3.6cm}}
\toprule
Method & Category & Mechanism & Extra inference cost & Metric families reported in original work & Limitation relevant here \\
\midrule
Boundary loss (Kervadec et al.) & loss & probability $\times$ signed distance, rebalanced with a region loss & none & overlap, HD & schedule-sensitive; backbone differs across papers \\
HD loss (Karimi \& Salcudean) & loss & distance-transform estimate of Hausdorff distance & none & overlap, HD, ASSD & computationally heavier distance maps \\
GSL (Celaya et al.) & loss & normalized distance-weighted error, region-to-surface annealing & none & overlap, HD95, ASSD & 3D origin; CPU distance transforms \\
Boundary DoU (Sun et al.) & loss & band-restricted difference-over-union & none & overlap, HD, Boundary IoU & single protocol per paper \\
Active boundary loss (Wang et al.) & loss & differentiable boundary alignment & none & boundary F-score, mIoU & natural-image evaluation \\
Psi-Net (Murugesan et al.) & multitask & mask, contour, distance heads & retained heads & overlap, HD & extra heads at inference \\
Gated-SCNN (Takikawa et al.) & architecture & gated shape stream & retained stream & mIoU, boundary F-score & heavy backbone \\
LB-UNet (Xu \& Tong) & lightweight architecture & boundary and prediction-map auxiliary modules & retained modules & overlap & binary skin-lesion only \\
EGE-UNet (Ruan et al.) & lightweight architecture & grouped multi-axis attention & retained & overlap & overlap-only evaluation \\
STDC (Fan et al.) & training-only head & detail head with boundary-derived targets, discarded at inference & none & mIoU, speed & natural-image evaluation \\
BPKD (Liu et al.) & distillation & boundary and body knowledge from a separate teacher & none & mIoU & requires a teacher network \\
MISSU; dual self-distillation (Wang et al.; Banerjee et al.) & self-distillation & distillation within one network & none & overlap, HD & 3D models \\
BDS-Lite (this work) & training recipe & boundary decoder $\to$ stop-gradient distillation + bounded gate; SDF loss & projection + gate (1,072 params) & overlap, boundary, distance, volume & compact 2D backbone only \\
\bottomrule
\end{tabular}}
\caption{Boundary-aware approaches. Metric families were checked against each source paper.}
\end{table}

\section*{S2. Training step and comparator objectives}
\begin{algorithm}[ht]
\caption{BDS-Lite training step and deployment}
\begin{algorithmic}[1]
\Require batch $(x,y)$, targets $y_b$ (boundary) and $\phi$ (SDF) precomputed from $y$
\State $\{f_\ell\}\gets\mathrm{Enc}(x)$;\quad $F_s\gets D_s(\{f_\ell\})$;\quad $q_s\gets P(F_s)$
\State $g\gets\sigma(W_gq_s+b_g)$;\quad $z_s\gets W_o\bigl(F_s\odot[1+\alpha(g-0.5)]\bigr)$
\State $(z_b,q_b)\gets D_b(\{f_\ell\})$ \Comment{training only}
\State $L\gets L_{\mathrm{ce}}(z_s,y)+L_{\mathrm{dice}}(z_s,y)+\lambda_{\mathrm{sdf}}L_{\mathrm{sdf}}(z_s,\phi)+\lambda_bL_{\mathrm{boundary}}(z_b,y_b)+\lambda_d\,\mathrm{SmoothL1}(q_s,\sg(q_b))$
\State backpropagate $L$ (mixed precision), clip gradient norm to 1, AdamW step
\State \textbf{Deployment:} delete $D_b$; keep Enc, $D_s$, $P$, gate, $W_o$; predict $\arg\max_k z_s$
\end{algorithmic}
\end{algorithm}

\paragraph{Soft Dice.} With $\varepsilon=10^{-6}$, $n$ indexing images in the batch and $u$ pixels,
\[
L_{\mathrm{dice}}=1-\frac{1}{K-1}\sum_{k=1}^{K-1}\frac{2\sum_{n,u}p_{nk}(u)y_{nk}(u)+\varepsilon}{\sum_{n,u}\left(p_{nk}(u)+y_{nk}(u)\right)+\varepsilon}.
\]

\paragraph{U-Net+GSL.} $L_{\mathrm{GSL}}=\alpha_tL_{\mathrm{reg}}+(1-\alpha_t)L_{\mathrm{surf}}$ is added to $L_{\mathrm{seg}}$ with $\lambda_{\mathrm{gsl}}=1$, where $L_{\mathrm{reg}}$ is cross-entropy plus squared-denominator Dice and
\[
L_{\mathrm{surf}}=1-\frac1N\sum_n\frac{\sum_kw'_{nk}\sum_u\left(D_{nk}(u)\,(1-y_{nk}(u)-p_{nk}(u))\right)^2}{\sum_kw'_{nk}\sum_uD_{nk}(u)^2},
\]
with $D_{nk}$ the exact Euclidean distance to class $k$ (image diagonal if absent) and $w'_{nk}=1/(1+(\sum_uy_{nk})^2)$. The schedule is the published step schedule: $\alpha_t=1-\lfloor(t-1)/5\rfloor/29$ for epochs $t\le145$ and $\alpha_t=0$ afterwards.

\paragraph{U-Net+Boundary Loss.} Following the official implementation, $L=w_r(t)L_{\mathrm{seg}}+w_b(t)\frac{1}{N(K-1)HW}\sum_{n,k,u}p_{nk}(u)\tilde\phi_{nk}(u)$ with unnormalized signed distances $\tilde\phi$ and the published rebalancing $w_r(1)=1$, $w_b(1)=0.01$, $w_r\leftarrow\max(0.1,w_r-0.01)$, $w_b\leftarrow w_b+0.01$ after every epoch.

\section*{S3. Datasets, arms, and comparator implementations}
{\small\input{generated/tab_descriptors.tex}}

\medskip\noindent Structural descriptors of the training masks (per foreground class). Empty frac.: fraction of images without the class; area fraction and number of 4-connected components per image; $L/A$: perimeter-to-area ratio for non-empty masks; $\partial$DSC/1\,px: $L/(2A)$. ISIC2018 masks have a median of \R{CMP.isic2018.median} component and \R{CMP.isic2018.single}\% have exactly one; the higher mean comes from a minority of fragmented masks. Synapse classes are numeric IDs; by area, emptiness and difficulty they correspond to aorta, gallbladder, left kidney, right kidney, liver, pancreas, spleen, and stomach.

\begin{table}[ht]\centering\small
\begin{tabular}{lcccc}
\toprule
Arm & Boundary decoder & Surface loss & Distillation & Gate \\
\midrule
U-Net & & & & \\
U-Net + surface loss & & yes & & \\
U-Net + boundary decoder & yes & & & \\
U-Net + decoder + surface loss & yes & yes & & \\
BDS-Lite $-$ surface loss & yes & & yes & yes \\
BDS-Lite $-$ distillation & yes & yes & & yes \\
BDS-Lite $-$ gate & yes & yes & yes & \\
BDS-Lite & yes & yes & yes & yes \\
\bottomrule
\end{tabular}
\caption{Components of each ablation arm. The first four rows form a two-by-two factorial of decoder and surface loss; the fourth was added under the frozen protocol (S13).}
\end{table}

\begin{table}[ht]\centering\small
\begin{tabular}{p{2.6cm}p{4.5cm}p{8.2cm}}
\toprule
Method & Source & Adaptation to the matched protocol \\
\midrule
U-Net+GSL & Celaya et al.; MIST GenSurfLoss (Apache-2.0) & 2D multi-class port; exact SciPy Euclidean distance transform; step schedule (5 epochs); $\lambda_{\mathrm{gsl}}=1$ added to CE+Dice. Retrained because the checkpoints behind the original rows are not available to us; the retrain reproduces the original training logs (S7). \\
U-Net+Boundary Loss & LIVIAETS/boundary-loss, commit 171c32d (MIT) & SurfaceLoss and one\_hot2dist as released; rebalancing schedule of the released \texttt{gdl\_surface\_steal} configuration; region term = CE+Dice of our U-Net instead of generalized Dice; all foreground classes. \\
Attention U-Net & in-repository implementation, width $c=32$ & identical encoder; additive attention gates on skip connections. \\
LB-UNet & xuxuxuxuxuxjh/LB-UNet, commit 913ee68 (no license; fetched at run time, not redistributed) & official architecture and widths; $K$-class output trained with CE+Dice; official deep-supervision weights on the foreground union and our union boundary target (the official genetic-algorithm key-point maps are binary-lesion specific); locked schedule instead of 256\,px/300 epochs/augmentation. Parameter count in our environment: 55,874 (2 classes) versus 0.038\,M reported in the paper. \\
\bottomrule
\end{tabular}
\caption{Comparator sources and adaptations. Configurations: \texttt{configs/revision/}; protocol: \texttt{revision/COMPARATOR\_PROTOCOL.md}.}
\end{table}

\clearpage
\section*{S4. Correction of the ISIC2018 evaluation}
\paragraph{Evidence.} The original ISIC2018 U-Net and BDS-Lite checkpoints record 226.9 optimizer steps per epoch (all six runs), equal to $1{,}815/8$ and not to $2{,}075/8$; ACDC (85) and Synapse (109) match their documented partitions. The best validation DSC in each ISIC training log equals, to six decimals, the evaluation on the 389-image validation set of the earlier split-seed-1 partition. Comparing the split manifests directly, of the 519 images of the documented validation set used for the original ISIC2018 results, 358 were in that partition's training set, 84 in its validation set, and 77 in its test set; the two partitions cover the same 2,594-image pool, and the official test set has no image identifier in common with it. Rerunning the checkpoints on the 519 images reproduces the original per-image values (maximum absolute DSC difference $4\times10^{-4}$ for U-Net and $1.1\times10^{-3}$ for BDS-Lite). All 18 original checkpoint files are byte-identical (SHA-256) to the archive manifest of the public release. With the original aggregation, the present evaluator reproduces every ACDC and Synapse value of the original Table~2 (per seed, maximum absolute difference $2\times10^{-4}$ for DSC, IoU and Boundary F1 and 0.03\,px for HD95 and ASSD) and of the original Table~3 (mean differences within 0.015; Wilcoxon statistics identical in 9 of 10 rows), with residual differences explained by CPU versus GPU arithmetic (\texttt{revision/validation/submitted\_tables\_reproduction.csv}, \texttt{checkpoint\_identity.json}; audit: \texttt{revision/GITHUB\_AUDIT.md}).

\paragraph{Originally reported values.} These values are preserved with their status; none was overwritten.
\begin{table}[ht]\centering
\resizebox{\linewidth}{!}{\input{generated/tab_s_original.tex}}
\caption{Originally submitted Table~2 (three-seed means; ISIC2018 on the 519-image set; original aggregation, in which classes absent from both prediction and reference count as perfect). Source: public release \texttt{results/main\_test\_means.csv} and \texttt{results/matched\_gsl\_summary.csv}.}
\end{table}

\paragraph{Repeated lesions.} Image IDs of the official test set do not occur in the development pool, but the ISIC archive contains repeated images of the same lesion under different IDs (re-captures, color or orientation changes). Each test image was compared with all 2,594 development images by the maximum Pearson correlation of channel-standardized $32\times32$ thumbnails over the eight flips and rotations (\texttt{analysis/revision/isic\_duplicate\_sensitivity.py}; per-image results in \texttt{revision/validation/isic\_test\_duplicates.csv}). Visual inspection showed certain repeats above $r=0.99$ and both repeats and distinct lesions between 0.95 and 0.97; neither correlation nor reference-mask agreement separates them cleanly. The contrasts were therefore repeated after excluding test images above four thresholds ($r>0.99$: \R{DUP.n.0.99} images; $>0.96$: \R{DUP.n.0.96}; $>0.93$: \R{DUP.n.0.93}; $>0.90$: \R{DUP.n.0.9}). The follow-up family equals the retrained set on ISIC2018.

{\small\input{generated/rl_tab_s_duplicates.tex}}

\clearpage
\section*{S5. BDS-Lite versus U-Net: full statistics and run-level variation}
\begin{table}[ht]\centering
\resizebox{\linewidth}{!}{\input{generated/rl_tab_primary_stats.tex}}
\caption{Original models. Oriented $\Delta$ is positive when BDS-Lite is better (distance differences negated). 95\% percentile bootstrap CI over units; Holm correction over five metrics within dataset. $^{*}$Holm $p<0.05$ with a CI excluding zero. SESOI: 0.01 (DSC, IoU, BF1), 2\,px (HD95), 0.5\,px (ASSD).}
\end{table}
\begin{table}[ht]\centering
\resizebox{\linewidth}{!}{\input{generated/rl_tab_s_matched.tex}}
\caption{Retrained models (same conventions).}
\end{table}

\paragraph{Run-level variation.} The margin is defined in the main text and was added after the results were known. ``Same-recipe gap'' is the difference between the three-run means of the original and retrained sets for one recipe, oriented so that positive favors the original set. Run means are weighted by test unit, so they differ slightly from the seed-averaged paired differences of the main tables; the two aggregations treat a class that is undefined in only some seeds differently. Per-run values: \texttt{revision/tables/families/run\_level\_noise.csv}.
\begin{table}[ht]\centering
\resizebox{\linewidth}{!}{\input{generated/tab_s_run_level.tex}}
\caption{Run-level variation and BDS-Lite differences (differences of run means).}
\end{table}

\begin{table}[ht]\centering\small
\input{generated/tab_s_selected_epochs.tex}
\caption{Selected checkpoint epochs (seeds 1/2/3) under the maximum-validation-DSC rule. Validation sets are small on ACDC (10 patients) and Synapse (4 cases), and validation curves fluctuate between epochs. The retrained set tends to select earlier epochs. The trainer of the original set is not preserved, so the cause cannot be established.}
\end{table}

\paragraph{ACDC U-Net gap.} Of the 9 other retrained arms, excluding LB-UNet, \R{RL.acdc.dsc.above.S0} reached the unit-weighted DSC of the original U-Net (\R{RL.acdc.dsc.S0.mean}). The two sets also selected different checkpoints: for U-Net, epochs \R{SEL.acdc.S0} in the original set and \R{SEL.acdc.M0} in the retrained set, from noisy validation curves on 10 patients.

\paragraph{Aggregation on Synapse.} Synapse slices differ in which organs they contain, so a value averaged first over samples and one averaged first over cases can differ. Retrained contrasts on Synapse DSC moved by up to 0.04 between such aggregations for one ablation (BDS-Lite without the surface loss), and we therefore treat small Synapse differences with caution. The main tables use the registered unit-weighted aggregation.

\begin{figure}[ht]\centering
\includegraphics[width=0.95\linewidth]{../figures/run_level.pdf}
\caption{Individual three-seed training runs on the test data. Markers: three runs per arm; vertical ticks: run means, weighted by test unit. Gray band: retrained U-Net mean $\pm$ run-level margin. The HD95 axes are reversed, so better values lie to the right in every panel. Arrows give the mean of arms outside the plotted range.}
\end{figure}

\paragraph{Follow-up family with class weights.} The retrained U-Net and BDS-Lite on ISIC2018 and Synapse come from a follow-up study that used the released pipeline; the training adapter used for all other retrained arms reproduces an archived run of that study exactly (identical selected epoch and per-image test metrics). In that study the ACDC BDS-Lite models used foreground class weights $(1.5,1.0,1.25)$, so we trained a further unweighted ACDC BDS-Lite set as the retrained reference. The follow-up family below therefore equals the retrained set on ISIC2018 and Synapse and differs on ACDC BDS-Lite only.
\begin{table}[ht]\centering
\resizebox{\linewidth}{!}{\input{generated/rl_tab_s_replication_means.tex}}
\caption{Means of the follow-up family (R) and the retrained references (M).}
\end{table}
\begin{table}[ht]\centering
\resizebox{\linewidth}{!}{\input{generated/rl_tab_s_replication.tex}}
\caption{Paired BDS-Lite versus U-Net analysis in the follow-up family.}
\end{table}

\clearpage
\section*{S6. Ablations}
\begin{table}[ht]\centering
\resizebox{\linewidth}{!}{\input{generated/tab_ablation_means.tex}}
\caption{Three-seed means $\pm$ seed SD of all ablation arms. ACDC decoder-only values use the corrected unweighted runs (S13). The image-level contrasts below are descriptive when comparing recipes (S14).}
\end{table}
\begin{table}[ht]\centering
\resizebox{\linewidth}{!}{\input{generated/tab_ablation_vs_full.tex}}
\caption{Each ablation relative to full BDS-Lite (negative: the removed component helped). $^{*}$Holm $p<0.05$ over all ablation--metric tests within dataset, with CI excluding zero.}
\end{table}
\begin{table}[ht]\centering
\resizebox{\linewidth}{!}{\input{generated/tab_ablation_vs_unet.tex}}
\caption{Each ablation relative to retrained U-Net; $^{*}$Holm $p<0.05$ with CI excluding zero.}
\end{table}

\clearpage
\section*{S7. Comparators, omnibus test, GSL reproduction, and published results}
All comparators were trained by the same adapter on the documented partitions; ISIC2018 values are on the official test set.
\begin{table}[ht]\centering
\resizebox{\linewidth}{!}{\input{generated/tab_comparator_means.tex}}
\caption{Comparator and retrained-reference means (three seeds; mean $\pm$ seed SD).}
\end{table}
\begin{table}[ht]\centering
\resizebox{\linewidth}{!}{\input{generated/tab_comparators.tex}}
\caption{Oriented mean differences of comparators against U-Net (upper panel) and full BDS-Lite (lower panel); $^{*}$Holm $p<0.05$ within family and dataset with CI excluding zero.}
\end{table}
\begin{table}[ht]\centering
\resizebox{\linewidth}{!}{\input{generated/rl_tab_s_comparators_vs_unet.tex}}
\caption{Comparators versus retrained U-Net. Holm correction over all comparator--metric tests within dataset.}
\end{table}
\begin{table}[ht]\centering
\resizebox{\linewidth}{!}{\input{generated/rl_tab_s_comparators_vs_full.tex}}
\caption{Comparators versus retrained BDS-Lite; positive values favor the comparator.}
\end{table}

\paragraph{Omnibus test.} Friedman test over the cluster means of all eleven retrained arms per dataset and metric ($\chi^2$ from \R{OM.min.chi2} to \R{OM.max.chi2}, $p\le$\,\R{OM.max.p}); mean ranks with 1 = best, lowest in bold. LB-UNet ranked last in \R{OM.acdc.C4.last} of 5 ACDC and \R{OM.synapse.C4.last} of 5 Synapse metrics, so these rejections partly reflect that arm. Full BDS-Lite never had the lowest mean rank; its positions ranged from \R{OM.M6.bestpos} to \R{OM.M6.worstpos} of \R{OM.arms}. Each arm contributes its own cluster means over the samples on which its metric is defined, so pairwise post-hoc differences are not identical to the paired family deltas. Post-hoc Wilcoxon tests with Holm correction over the 55 pairs per dataset and metric: \texttt{revision/tables/families/omnibus\_posthoc.csv}.
\begin{table}[ht]\centering
\resizebox{\linewidth}{!}{\input{generated/tab_s_omnibus.tex}}
\caption{Friedman $\chi^2$, $p$ and mean ranks of all retrained arms.}
\end{table}

\paragraph{Reproduction of the original U-Net+GSL.} The checkpoints behind the originally submitted U-Net+GSL rows are not available to us, so U-Net+GSL was retrained with the original configuration and the same loss implementation on the documented partitions. Each retrained run was compared epoch by epoch with the original training log (\texttt{analysis/revision/gsl\_trajectory\_compare.py}). In \R{GSLTR.equal} of \R{GSLTR.complete} runs (\R{GSLTR.datasets}) the validation DSC equals the original value at every epoch (maximum absolute difference \R{GSLTR.maxdiff}), so the same checkpoint epoch is selected. Because the retrain reads the documented partitions through exact manifests, this also establishes that the original U-Net+GSL ISIC2018 models were trained on the documented 2,075-image set, unlike the original U-Net and BDS-Lite models. Their originally reported ISIC2018 value was the best-validation value on the 519-image selection set. With the original aggregation, the \R{GSLT2.runs} retrained runs reproduce the originally submitted U-Net+GSL values (maximum absolute difference \R{GSLT2.overlap} for DSC, IoU and Boundary F1 and \R{GSLT2.distance}\,px for HD95 and ASSD).
\begin{table}[ht]\centering
\resizebox{\linewidth}{!}{\input{generated/tab_s_gsl_trajectory.tex}}
\caption{Retrained U-Net+GSL versus the original training logs (validation DSC as logged during training).}
\end{table}

\paragraph{Published results.} Every value was verified against the cited page or table (\texttt{revision/lit\_numbers.csv}). Not directly comparable with our study: protocols differ in training cases (Synapse: 18 versus our 14), pretraining, augmentation, resolution, split (ISIC2018 lightweight models use a random 7:3 split), and metric aggregation.
\begin{table}[ht]\centering\scriptsize
\resizebox{\linewidth}{!}{\input{generated/tab_literature.tex}}
\caption{Published results under standard protocols.}
\end{table}

\clearpage
\section*{S8. Error analysis}
\paragraph{All-case distance sensitivity.} Primary HD95 and ASSD omit one-sided empty pairs. The registered all-case analysis assigns them the image diagonal (316.8\,px) and assigns zero to both-empty pairs. It is a failure penalty, not a measured distance. Positive values favor BDS-Lite.
\begin{table}[ht]\centering\small
\input{generated/tab_s_all_case.tex}
\caption{Primary and all-case distance differences, BDS-Lite versus U-Net. Source: \texttt{revision/tables/families/all\_case\_sensitivity.csv}.}
\end{table}

\paragraph{Error location.} Pixels or components per image, seed means. Near: within 2\,px of the reference contour. Spurious component: a predicted component with no reference overlap. Spurious components per image rose under BDS-Lite in both three-seed model sets: from \R{EB.isic2018.S0.spurious} to \R{EB.isic2018.S6.spurious} and \R{EB.isic2018.M0.spurious} to \R{EB.isic2018.M6.spurious} on ISIC2018, \R{EB.acdc.S0.spurious} to \R{EB.acdc.S6.spurious} and \R{EB.acdc.M0.spurious} to \R{EB.acdc.M6.spurious} on ACDC, and \R{EB.synapse.S0.spurious} to \R{EB.synapse.S6.spurious} and \R{EB.synapse.M0.spurious} to \R{EB.synapse.M6.spurious} on Synapse.
{\small\input{generated/rl_tab_s_errorbands.tex}}

\paragraph{Missed structures and false alarms by reference area.}
{\small\input{generated/rl_tab_s_collapse.tex}}

\paragraph{Per-class DSC.}
{\small\input{generated/rl_tab_s_perclass.tex}}

\clearpage
\section*{S9. Volume-level evaluation}
Slice predictions were stacked into volumes. Distances are in index units (ACDC processed 224 grid; Synapse original H5 grid), not millimeters. Absolute 3D DSC was \R{VOL.acdc.S0.dsc}\% and \R{VOL.acdc.S6.dsc}\% for the original U-Net and BDS-Lite on ACDC, and \R{VOL.synapse.S0.dsc}\% and \R{VOL.synapse.S6.dsc}\% on Synapse. Volume metrics penalize slice-wise inconsistency, which 2D training does not constrain. The follow-up family equals the retrained set on Synapse; on ACDC its BDS-Lite is class-weighted (S5).
\begin{table}[ht]\centering
\resizebox{0.8\linewidth}{!}{\input{generated/rl_tab_volume.tex}}
\caption{Volume-level BDS-Lite versus U-Net differences per patient (ACDC, frames averaged) or case (Synapse), oriented so that positive values favor BDS-Lite. Holm correction over the three endpoints within dataset and family.}
\end{table}

\clearpage
\section*{S10. Gradient diagnostics}
In seed-1 diagnostic runs of the full recipe we computed, every 50 steps, the gradients of the segmentation loss and of each auxiliary loss on the shared encoder parameters. These runs are single-seed, and the ACDC run used the class-weighted configuration, so they are indicative only. The cosine between the auxiliary and segmentation gradients had means of \R{GR.isic2018.aux.all.cos}, \R{GR.acdc.aux.all.cos}, and \R{GR.synapse.aux.all.cos} on ISIC2018, ACDC, and Synapse. Negative cosines were rare (\R{GR.isic2018.aux.all.conflict}\%, \R{GR.acdc.aux.all.conflict}\%, and \R{GR.synapse.aux.all.conflict}\% of probes), so we do not interpret these values as gradient conflict. With one run per dataset, the ordering across datasets is an observation, not a test.

{\small\input{generated/rl_tab_s_gradient.tex}}

\emph{Numerical consistency.} The auxiliary gradient and its three components were obtained in separate FP32 backward passes (cuDNN convolutions with TF32 enabled), so the linearity identity $|g_{\mathrm{aux}}|\cos_{\mathrm{aux}}=\sum_k|g_k|\cos_k$ holds only up to accumulation error, and exact equivalence of the separate and combined passes is not claimed. A strict check of this identity failed during development, and the failure log is retained. Over all probes of the three runs, the median relative residual was at most \R{GRA.rel.median.max}; attributing the whole residual to the auxiliary cosine changes it by at most \R{GRA.cos_err.max} (99th percentile \R{GRA.cos_err.p99}), far below the smallest difference between the dataset means (\R{GRA.gap.min}). The residual therefore cannot affect the reported cosines or their ordering (\texttt{analysis/revision/gradient\_residual\_audit.py}).

\begin{figure}[ht]\centering
\includegraphics[width=\linewidth]{../figures/gradient_cosine.pdf}
\caption{Cosine similarity between the segmentation gradient and auxiliary gradients on shared encoder parameters, every 50 steps (dots) with 25-probe moving average (lines), seed-1 diagnostic runs.}
\end{figure}

\clearpage
\section*{S11. Sensitivity to $\alpha$ and $\lambda$}
The gate scale was varied over $\alpha\in\{0,0.1,0.25,0.5,1.0\}$ with three seeds on ISIC2018 and one seed on Synapse, and the auxiliary weights over $\lambda_b=\lambda_d\in\{0.01,0.05,0.2\}$ with one seed on ISIC2018 (\texttt{revision/SENSITIVITY\_PROTOCOL.md}; Synapse $\alpha$ and ISIC2018 $\lambda$ were trimmed to one seed before training). Single-seed rows are descriptive, and no value was selected from these runs.

\emph{Gate scale, ISIC2018.} Against $\alpha=0.25$, none of the \R{SENS.isic2018.alpha.M6.tests} contrasts lay beyond the threshold (\R{SENS.isic2018.alpha.M6.equivalent} within it, the others inconclusive). Against U-Net, \R{SENS.isic2018.alpha.M0.equivalent} of the \R{SENS.isic2018.alpha.M0.tests} contrasts were within the threshold. The response was not monotonic. Validation DSC ranged from \R{SENSM.isic2018.alpha.val.min} to \R{SENSM.isic2018.alpha.val.max} (highest at $\alpha=$\,\R{SENSM.isic2018.alpha.val.argmax}; registered value \R{SENSM.isic2018.alpha.0.25.val}), so $\alpha=0.25$ is an a-priori choice that the validation data do not single out.

\emph{Gate scale, Synapse (seed 1).} Overlap intervals were inconclusive and HD95 stayed within the threshold for every value (for example DSC \R{SENS.synapse.alpha.0.M6.dsc.delta}, \R{SENS.synapse.alpha.0.M6.dsc.ci}, for $\alpha=0$). Every $\alpha$, including $\alpha=0$, had higher DSC than the seed-1 U-Net, which reflects that U-Net run, the weakest of its three seeds (DSC \R{SEED.synapse.M0.dsc.s1} versus \R{SEED.synapse.M0.dsc.s2} and \R{SEED.synapse.M0.dsc.s3}).

\emph{Auxiliary weights, ISIC2018 (seed 1).} Both alternatives gave higher seed-1 test values than $\lambda=0.05$ on every endpoint (Boundary F1 \R{SENS.isic2018.lambda.0.01.M6.boundary_f1.delta} at $\lambda=0.01$ and \R{SENS.isic2018.lambda.0.2.M6.boundary_f1.delta} at $\lambda=0.2$), whereas validation DSC did not distinguish the three values (\R{SENSM.isic2018.lambda.0.01.val}, \R{SENSM.isic2018.lambda.0.05.val}, \R{SENSM.isic2018.lambda.0.2.val}). The $\lambda=0.05$ seed-1 run was the weakest of its three seeds on the test set (Boundary F1 \R{SEED.isic2018.M6.boundary_f1.s1} versus \R{SEED.isic2018.M6.boundary_f1.s2} and \R{SEED.isic2018.M6.boundary_f1.s3}), so part of this difference is run-to-run variation. A multi-seed study would be needed to separate the two.
\begin{table}[ht]\centering
\resizebox{\linewidth}{!}{\input{generated/tab_s_sensitivity_means.tex}}
\caption{Validation and test means per setting (reg. = registered value). ISIC2018 validation DSC is on the 519 documented validation images; ``--'' where the reference run's validation record is not on the same pipeline.}
\end{table}
\begin{table}[ht]\centering
\resizebox{\linewidth}{!}{\input{generated/tab_s_sensitivity_stats.tex}}
\caption{Oriented paired differences of each setting relative to the registered BDS-Lite value and to U-Net (same seeds), positive when the setting is better. $^{*}$Holm $p<0.05$ with CI excluding zero (three-seed ISIC2018 $\alpha$ family only); $^{\uparrow}$/$^{\downarrow}$ interval beyond $\pm$SESOI. Single-seed rows carry no Holm decision. Full intervals: \texttt{revision/tables/families/sensitivity.csv}.}
\end{table}

\clearpage
\section*{S12. Qualitative examples}
Selection and display rules are those of the main qualitative figure. ACDC and Synapse panels show label maps only (see the Datasets section of the main text). False-alarm classes, absent from the reference, do not enter the class-macro HD95 and are reported as false-alarm rates.
\begin{figure}[ht]\centering
\includegraphics[width=0.85\linewidth]{../figures/qualitative_acdc.pdf}
\caption{ACDC qualitative examples (label maps).}
\end{figure}
\begin{figure}[ht]\centering
\includegraphics[width=0.85\linewidth]{../figures/qualitative_synapse.pdf}
\caption{Synapse qualitative examples (label maps).}
\end{figure}


\clearpage
\section*{S13. Additional trainings: protocols, results, and mechanism}
\paragraph{Protocol.} The additional trainings are defined in \texttt{revision/EXPANSION\_PROTOCOL.md}. Before they finished, \texttt{revision/NARRATIVE\_CRITERIA\_ADDENDUM.md} fixed which results would justify a stronger claim for BDS-Lite, and none of its stronger tiers (1--3) was reached. Both files are in version control with commit times preceding the results they govern. Configurations, sources, and manifests are hash-checked by the queue runner. Rebuilt seed-1 configurations of the four ISIC2018 arms replayed the archived validation trajectory exactly over four epochs (\texttt{revision/expansion\_gate/gate.json}). All 40 runs passed the checks of \texttt{revision/validation/expansion\_validation.json} before any of their test metrics was read.

\paragraph{Correction of the ACDC decoder-only arm.} The first ACDC decoder-only runs used foreground class weights $(1.5,1.0,1.25)$, unlike every other retrained ACDC arm, contrary to the design. Three unweighted runs replaced them in all analyses. Only contrasts involving this arm changed, and no Holm decision changed. The earlier runs are retained.

\paragraph{Operational notes.} One run paused for 17 minutes after an unrelated process occupied the GPU during a cooling stop, and resumed from its saved epoch; recorded in \texttt{revision/DEVIATIONS.md}.

\begin{table}[ht]\centering
\resizebox{\linewidth}{!}{\input{generated/tab_s_ten_seed_all.tex}}
\caption{Frozen contrasts on ISIC2018, all metrics. Positive favors the first arm, with distance differences negated. Holm correction over six contrasts within each metric on fresh seeds 4--10.}
\end{table}

\begin{table}[ht]\centering
\resizebox{\linewidth}{!}{\input{generated/tab_s_interaction.tex}}
\caption{Factorial and transfer-path effects on Boundary F1, seeds 1--3, with 95\% intervals. Interaction: $(\text{decoder+surface}-\text{decoder})-(\text{surface}-\text{U-Net})$. Pooled by inverse variance. DSC and HD95 for the central contrasts are in the main text; all metrics: \texttt{revision/tables/families/expansion\_interaction.csv}.}
\end{table}

\begin{table}[ht]\centering
\resizebox{0.9\linewidth}{!}{\input{generated/tab_s_mechanism.tex}}
\caption{What distillation and the gate do. Inference only, post hoc, seeds 1--3. Average precision (AP) against the union boundary target on up to 300 test images per dataset; readouts are closed-form ridge fits on validation pixels. Distilled arms: BDS-Lite and its variants without the surface loss or the gate. Undistilled arms: decoder-only and BDS-Lite without distillation. Gate structure uses 96 test images per dataset. Scripts: \texttt{analysis/revision/boundary\_transfer\_probe.py}, \texttt{gate\_structure.py}.}
\end{table}

\clearpage
\section*{S14. Repeatability of image-level tests}
This analysis was designed after the ten-seed results were known and uses completed runs only (\texttt{analysis/revision/seed\_subset\_analysis.py}). For one recipe, two disjoint groups of $k$ of its ten runs are compared with the Wilcoxon signed-rank test on per-image differences of group means, as in the registered families. Any significant result is therefore a difference between trained models and not between recipes. At most 2,000 splits are sampled per recipe. The table reports how often $p<0.05$ occurred. Rates stay high as $k$ grows from 3 to 5 because the image-level test measures how consistently two fixed sets of models differ across images, and this does not shrink with averaging in the way a comparison of recipes would. Results for draws of two different arms are in \texttt{revision/tables/families/seed\_subset.csv}; rows with $k=7$ share most of their runs, so their sign agreement is high by construction and should not be interpreted.
\begin{table}[ht]\centering\small
\input{generated/tab_s_seed_subset.tex}
\caption{Same-recipe comparisons on the ISIC2018 official test set.}
\end{table}

\end{document}
